{% Q2388375[94F]
\needspace{4\baselineskip} 
\item \rtask \ponto{\pt} 
Em relação à Linguagem de Programação Python, julgue o item.

A função \lstinline[style=Python]|range()|, em Python, é empregada, especificamente, para adicionar um item ao final de uma lista existente.

% F
{\setlength{\columnsep}{0pt}\renewcommand{\columnseprule}{0pt}
\begin{multicols}{2}
\begin{answerlist}[label={\texttt{\Alph*}.},leftmargin=*]
    \ti[V.]
    \ifnum\gabarito=1\doneitem[F.]\else\ti[F.]\fi % gabarito
\end{answerlist}
\end{multicols}
}

% Gabarito: E - Errado
% 
% A função range() em Python não tem o propósito de adicionar itens ao final de uma lista. Pelo contrário, seu uso principal é gerar sequências numéricas, não modificar estruturas de dados como listas. Mais especificamente, range() é frequentemente utilizada em laços de repetição, como for, para iterar sobre uma sequência de números.
% 
% Por exemplo:
% 
% <span style="color: blue;">for i in range(0, 5):
%     print(i)
% # Isso imprimirá os números de 0 a 4</span>
% 
% Se o objetivo é adicionar um item ao final de uma lista em Python, usamos o método append() da lista. Por exemplo:
% 
% <span style="color: blue;">minha_lista = [1, 2, 3]
% minha_lista.append(4)
% print(minha_lista)
% # Isso resultará em [1, 2, 3, 4]</span>
% 
% Portanto, a afirmação da questão está incorreta, pois confunde a finalidade da função range() em Python, levando a uma interpretação equivocada de suas funcionalidades.
}

